%!TEX root = ./syllabus-cis4930-cis5930.tex

\subsection*{Course Evaluation and Grading Policies} 

\begin{outline} 
\1 \textbf{Assignments: 45\%}
\end{outline} 

There will be several homework assignments (written and coding-based) spaced out over the course of the semester. All assignments are equally weighted. Submission instructions will be posted on GitHub Classroom and/or Canvas; most project submissions use GitHub. \\

% \begin{outline} 
% \1 \textbf{In-Class Quizzes: 2\%}
% \end{outline} 

% There will be several short in-class quizzes throughout the semester. These quizzes are primarily used to track attendance and encourage active participation.


\begin{outline} 
\1 \textbf{Paper Practicum (Reading \& Presentation): 15\%}
\end{outline}

Each student will complete \textbf{one} in-class paper practicum on top-tier conference papers related to \textbf{Future Edge Networks}.

\begin{enumerate}[label=(\arabic*), leftmargin=2em]
%   \item \textbf{Requirement:} Complete \textbf{three} practicums during the semester (total \textbf{15\%}).
  
  \item \textbf{Format: Individual presentation only} (one person per practicum). The same time allocation applies to both undergraduate and graduate students: \textbf{15 minutes} presentation + \textbf{5 minutes} Q\&A.

  \item \textbf{Scheduling:} \textbf{Three} papers are presented per class meeting. The practicum block runs for roughly seven weeks after the introductory lectures conclude.

  \item \textbf{Paper Assignment:} Papers will be assigned by the instructor. Students will first submit their preferred presentation dates (and optionally topics of interest). The instructor will then assign a paper to each student based on the course schedule and topic sequence.
\end{enumerate}

\textbf{Important policy for paper practicum:}

\begin{enumerate}[label=(\arabic*), leftmargin=2em]
  \item \textbf{In-person presentation only.} No make-up or deferred practicums will be approved.
  \item Presentation time slots will be assigned in random order. Plan ahead when submitting your availability.
\end{enumerate}


\begin{outline} 
\1 \textbf{Final Project Report \& Presentation: 40\%}
\end{outline}


\begin{enumerate}[label=(\arabic*), leftmargin=2em]
	\item \textbf{Final presentation (15\%)} will take place at the end of the semester (dates subject to change). Each person will have \textbf{15 minutes} for presentation and \textbf{5 minutes} for Q\&A.

	\item \textbf{Demo (25\%)} You may choose whether or not to share the demo code. However, you must present a working demo during your project presentation.
\end{enumerate}






\textbf{Important policy for final presentation:}

\begin{enumerate}[label=(\arabic*), leftmargin=2em]
  \item \textbf{In-person presentation only.} No make-up or deferred presentations will be approved.
  \item Multiple time slots will be assigned for each class in random order among presenters. Plan ahead when submitting your availability.
\end{enumerate}

Please see a detailed introduction of the final project report and presentation on the course site.

\textbf{Extra Bonus.}

Students are encouraged to prepare submissions to arXiv or major AI/ML/DM conferences based on their projects. Please make an appointment with the instructor prior to any submission plan for a comprehensive evaluation of the research topic. Each submission under the instructor's recognition will gain \textbf{2 points} on the final grade.


Letter grades are assigned based on the final percent using the cut-off values:

\begin{flushleft}
	\textbf{Grade} 		\hfill 		\textbf{\% bound} \\
	\vspace{.1in} \hrule  \vspace{.1in}
	
		A 		\dotfill			93 - 100  \\
		A- 		\dotfill 			90 - 92.9 \\
		B+ 		\dotfill 			87 - 89.9 \\
		B 		\dotfill 			83 - 86.9 \\
		B- 		\dotfill 			80 - 82.9 \\
		C+ 		\dotfill 			77 - 79.9 \\
		C 		\dotfill 			73 - 76.9 \\
		C- 		\dotfill 			70 - 72.9 \\
		D 		\dotfill 			60 - 69.9 \\
		F 		\dotfill 			<60       \\
\end{flushleft}
